\documentclass{ximera}

\input{../../../preamble.tex}



\definecolor{fvdnavy}{HTML}{071D43}
\definecolor{fvdcyan}{HTML}{20BFC6}
\definecolor{fvdcyanlight}{HTML}{EFFCFB}
\definecolor{fvdmagenta}{HTML}{F10D69}
\definecolor{fvdmagentalight}{HTML}{FFF1F7}
\definecolor{fvdtext}{HTML}{163044}





%=================================================
% Pedagogische omgevingen
% PDF: tcolorbox
% HTML: learning-box
%=================================================

\newenvironment{voorkennis}
{%
    \ifdefined\HCode
        \HCode{%
<section class="learning-box learning-box--prior">
<div class="learning-box__header">
<span class="learning-box__icon" aria-hidden="true"></span>
<span class="learning-box__title">Voorkennis</span>
</div>
<div class="learning-box__content">
        }%
        \begin{itemize}
    \else
       \begin{tcolorbox}[
    enhanced,
    breakable,
    colback=fvdcyanlight,
    colframe=fvdcyan,
    coltext=fvdtext,
    boxrule=0.5pt,
    leftrule=4pt,
    arc=4mm,
    outer arc=4mm,
    left=7mm,
    right=6mm,
    top=5mm,
    bottom=5mm,
    before skip=6mm,
    after skip=6mm,
    title=Voorkennis,
    coltitle=fvdcyan,
    fonttitle=\bfseries\Large,
    attach title to upper={\par\smallskip},
    boxed title style={
        empty,
        left=0mm,
        right=0mm,
        top=0mm,
        bottom=0mm
    },
    overlay={
        \draw[
            fvdcyan,
            line width=0.5pt
        ]
        ([xshift=42mm,yshift=-13mm]frame.north west)
        --
        ([xshift=-7mm,yshift=-13mm]frame.north east);
    }
]
        \begin{itemize}
    \fi
}
{%
    \end{itemize}
    \ifdefined\HCode
        \HCode{%
</div>
</section>
        }%
    \else
        \end{tcolorbox}
    \fi
}


\newenvironment{leerdoelen}
{%
    \ifdefined\HCode
        \HCode{%
<section class="learning-box learning-box--goals">
<div class="learning-box__header">
<span class="learning-box__icon" aria-hidden="true"></span>
<span class="learning-box__title">Leerdoelen</span>
</div>
<div class="learning-box__content">
        }%
        \begin{itemize}
    \else
\begin{tcolorbox}[
    enhanced,
    breakable,
    colback=fvdmagentalight,
    colframe=fvdmagenta,
    coltext=fvdtext,
    boxrule=0.5pt,
    leftrule=4pt,
    arc=4mm,
    outer arc=4mm,
    left=7mm,
    right=6mm,
    top=5mm,
    bottom=5mm,
    before skip=6mm,
    after skip=6mm,
    title=Leerdoelen,
    coltitle=fvdmagenta,
    fonttitle=\bfseries\Large,
    attach title to upper={\par\smallskip},
    boxed title style={
        empty,
        left=0mm,
        right=0mm,
        top=0mm,
        bottom=0mm
    },
    overlay={
        \draw[
            fvdmagenta,
            line width=0.5pt
        ]
        ([xshift=42mm,yshift=-13mm]frame.north west)
        --
        ([xshift=-7mm,yshift=-13mm]frame.north east);
    }
]
        \begin{itemize}
    \fi
}
{%
    \end{itemize}
    \ifdefined\HCode
        \HCode{%
</div>
</section>
        }%
    \else
        \end{tcolorbox}
    \fi
}



\addPrintStyle{../../..}

\title{Oefeningen --- Uitspraken verbinden}

\author{Ingmar Herreman}



\begin{document}

% FVD_CHAPTER_DATA_START
% Automatisch gegenereerd uit index.tex.
% Niet handmatig aanpassen.
\fvdchapterdata{1}{Logica — De taal van correct redeneren}{2}
% FVD_CHAPTER_DATA_END

















\maketitle





% ============================================================

% BASIS

% ============================================================



\begin{oefening}[De negatie bepalen]{\oefB \ESS}

Geef telkens de negatie van de uitspraak.

\begin{question}

\[
p:\quad 7<12
\]

\begin{oplossing}

De negatie is

\[
\neg p:\quad 7\geq12.
\]

\end{oplossing}

\end{question}



\begin{question}

\[
q:\quad x>5
\]

\begin{oplossing}

De negatie is

\[
\neg q:\quad x\leq5.
\]

Let op: $x<5$ is niet de volledige negatie, want dan wordt $x=5$

vergeten.

\end{oplossing}

\end{question}



\begin{question}

\[
r:\quad x\leq-2
\]

\begin{oplossing}

\[
\neg r:\quad x>-2.
\]

\end{oplossing}

\end{question}



\begin{question}

\[
s:\quad n=8
\]

\begin{oplossing}

\[
\neg s:\quad n\neq8.
\]

\end{oplossing}

\end{question}



\begin{question}

\[
t:\quad n\text{ is deelbaar door }3.
\]

\begin{oplossing}

\[
\neg t:\quad n\text{ is niet deelbaar door }3.
\]

\end{oplossing}

\end{question}

\end{oefening}





\begin{oefening}[Waarheidswaarde van een negatie]{\oefB \ESS}

Gegeven zijn de uitspraken

\[
p:\quad 13\text{ is een priemgetal},
\]

\[
q:\quad 20<15,
\]

\[
r:\quad \sqrt{16}=4.
\]

Bepaal telkens de waarheidswaarde.

\begin{question}

$p$

\begin{oplossing}

Waar.

\end{oplossing}

\end{question}



\begin{question}

$\neg p$

\begin{oplossing}

Onwaar.

\end{oplossing}

\end{question}



\begin{question}

$q$

\begin{oplossing}

Onwaar.

\end{oplossing}

\end{question}



\begin{question}

$\neg q$

\begin{oplossing}

Waar.

\end{oplossing}

\end{question}



\begin{question}

$\neg r$

\begin{oplossing}

$r$ is waar, want

\[
\sqrt{16}=4.
\]

Dus $\neg r$ is onwaar.

\end{oplossing}

\end{question}

\end{oefening}





\begin{oefening}[En: de conjunctie]{\oefB \ESS}

Gegeven zijn

\[
p:\quad 18\text{ is even}
\]

en

\[
q:\quad 18\text{ is deelbaar door }3.
\]

\begin{question}

Bepaal de waarheidswaarde van $p$ en $q$.

\begin{oplossing}

Beide uitspraken zijn waar.

\[
p=W,
\qquad
q=W.
\]

\end{oplossing}

\end{question}



\begin{question}

Bepaal de waarheidswaarde van

\[
p\land q.
\]

\begin{oplossing}

Een conjunctie is waar als beide onderdelen waar zijn.

Dus

\[
p\land q=W.
\]

\end{oplossing}

\end{question}



\begin{question}

Schrijf $p\land q$ in woorden.

\begin{oplossing}

\begin{center}

$18$ is even en $18$ is deelbaar door $3$.

\end{center}

\end{oplossing}

\end{question}

\end{oefening}





\begin{oefening}[Of: de disjunctie]{\oefB \ESS}

Gegeven zijn

\[
p:\quad 15\text{ is even}
\]

en

\[
q:\quad 15\text{ is deelbaar door }5.
\]

\begin{question}

Bepaal de waarheidswaarde van $p$ en $q$.

\begin{oplossing}

\[
p=O,
\qquad
q=W.
\]

\end{oplossing}

\end{question}



\begin{question}

Bepaal de waarheidswaarde van

\[
p\lor q.
\]

\begin{oplossing}

De disjunctie is waar zodra minstens één van beide uitspraken waar is.

Omdat $q$ waar is,

\[
p\lor q=W.
\]

\end{oplossing}

\end{question}



\begin{question}

Bepaal de waarheidswaarde van

\[
p\land q.
\]

\begin{oplossing}

Omdat $p$ onwaar is, zijn niet beide uitspraken waar.

Dus

\[
p\land q=O.
\]

\end{oplossing}

\end{question}

\end{oefening}





\begin{oefening}[Conjunctie of disjunctie?]{\oefB \ESS}

Gegeven zijn

\[
p=W,
\qquad
q=O.
\]

Bepaal de waarheidswaarde van elke uitdrukking.

\begin{question}

\[
p\land q
\]

\begin{multipleChoice}

  \choice{waar}

  \choice[correct]{onwaar}

\end{multipleChoice}

\end{question}



\begin{question}

\[
p\lor q
\]

\begin{multipleChoice}

  \choice[correct]{waar}

  \choice{onwaar}

\end{multipleChoice}

\end{question}



\begin{question}

\[
\neg p
\]

\begin{multipleChoice}

  \choice{waar}

  \choice[correct]{onwaar}

\end{multipleChoice}

\end{question}



\begin{question}

\[
\neg q
\]

\begin{multipleChoice}

  \choice[correct]{waar}

  \choice{onwaar}

\end{multipleChoice}

\end{question}

\end{oefening}





\begin{oefening}[Het wiskundige ``of'']{\oefB \ESS}

Gegeven zijn

\[
p:\quad 12\text{ is even}
\]

en

\[
q:\quad 12\text{ is deelbaar door }3.
\]

\begin{question}

Welke uitspraken zijn waar?

\begin{multipleChoice}

  \choice{$p$ is waar en $q$ is onwaar.}

  \choice{$p$ is onwaar en $q$ is waar.}

  \choice[correct]{$p$ en $q$ zijn beide waar.}

  \choice{$p$ en $q$ zijn beide onwaar.}

\end{multipleChoice}

\end{question}



\begin{question}

Wat is de waarheidswaarde van

\[
p\lor q?
\]

\begin{multipleChoice}

  \choice[correct]{waar}

  \choice{onwaar}

\end{multipleChoice}

\end{question}



\begin{question}

Een leerling zegt:

\begin{quote}

``$p\lor q$ kan niet waar zijn, want bij ``of'' mag maar één van de

twee uitspraken waar zijn.''

\end{quote}

Beoordeel deze redenering.

\begin{oplossing}

De redenering is fout.

In de wiskundige logica is $\lor$ een \textbf{inclusieve disjunctie}.

Dat betekent dat

\[
p\lor q
\]

waar is wanneer minstens één van beide uitspraken waar is.

Ook wanneer $p$ en $q$ allebei waar zijn, is $p\lor q$ dus waar.

\end{oplossing}

\end{question}

\end{oefening}





\begin{oefening}[Waarheidstabel van de basisoperatoren]{\oefB \ESS}

Vul de waarheidstabel aan.

\[
\begin{array}{c|c|c|c|c}
p&q&\neg p&p\land q&p\lor q\\
\hline
W&W&\phantom{O}&\phantom{W}&\phantom{W}\\
W&O&\phantom{O}&\phantom{O}&\phantom{W}\\
O&W&\phantom{W}&\phantom{O}&\phantom{W}\\
O&O&\phantom{W}&\phantom{O}&\phantom{O}
\end{array}
\]

\begin{oplossing}

\[
\begin{array}{c|c|c|c|c}
p&q&\neg p&p\land q&p\lor q\\
\hline
W&W&O&W&W\\
W&O&O&O&W\\
O&W&W&O&W\\
O&O&W&O&O
\end{array}
\]

\end{oplossing}

\end{oefening}





\begin{oefening}[Van symbolen naar woorden]{\oefB \ESS}

Gegeven zijn

\[
p:\quad x>2
\]

en

\[
q:\quad x<7.
\]

Vertaal elke uitdrukking naar gewone wiskundige taal.

\begin{question}

\[
p\land q
\]

\begin{oplossing}

\[
x>2\text{ en }x<7.
\]

Dit kan ook worden geschreven als

\[
2<x<7.
\]

\end{oplossing}

\end{question}



\begin{question}

\[
p\lor q
\]

\begin{oplossing}

\[
x>2\text{ of }x<7.
\]

Het woord ``of'' wordt inclusief geïnterpreteerd.

\end{oplossing}

\end{question}



\begin{question}

\[
\neg p
\]

\begin{oplossing}

\[
x\leq2.
\]

\end{oplossing}

\end{question}



\begin{question}

\[
\neg q
\]

\begin{oplossing}

\[
x\geq7.
\]

\end{oplossing}

\end{question}

\end{oefening}





\begin{oefening}[Van woorden naar symbolen]{\oefB \ESS}

Gegeven zijn

\[
p:\quad n\text{ is even}
\]

en

\[
q:\quad n>10.
\]

Schrijf de volgende uitspraken met logische symbolen.

\begin{question}

$n$ is even en groter dan $10$.

\begin{oplossing}

\[
p\land q.
\]

\end{oplossing}

\end{question}



\begin{question}

$n$ is even of groter dan $10$.

\begin{oplossing}

\[
p\lor q.
\]

\end{oplossing}

\end{question}



\begin{question}

$n$ is niet even.

\begin{oplossing}

\[
\neg p.
\]

\end{oplossing}

\end{question}



\begin{question}

$n$ is even en niet groter dan $10$.

\begin{oplossing}

\[
p\land\neg q.
\]

\end{oplossing}

\end{question}

\end{oefening}





\begin{oefening}[Samengestelde uitspraken berekenen]{\oefB \ESS}

Gegeven zijn

\[
p=W,
\qquad
q=O,
\qquad
r=W.
\]

Bepaal telkens de waarheidswaarde.

\begin{question}

\[
\neg p\lor q
\]

\begin{oplossing}

\[
\neg p=O.
\]

Dus

\[
\neg p\lor q
=
O\lor O
=
O.
\]

\end{oplossing}

\end{question}



\begin{question}

\[
p\land\neg q
\]

\begin{oplossing}

Omdat

\[
\neg q=W,
\]

krijgen we

\[
p\land\neg q
=
W\land W
=
W.
\]

\end{oplossing}

\end{question}



\begin{question}

\[
(p\land q)\lor r
\]

\begin{oplossing}

Eerst:

\[
p\land q
=
W\land O
=
O.
\]

Daarna:

\[
O\lor W=W.
\]

Dus de samengestelde uitspraak is waar.

\end{oplossing}

\end{question}



\begin{question}

\[
p\land(q\lor r)
\]

\begin{oplossing}

Eerst:

\[
q\lor r
=
O\lor W
=
W.
\]

Daarna:

\[
p\land W
=
W\land W
=
W.
\]

\end{oplossing}

\end{question}

\end{oefening}





% ============================================================

% VERDIEPING

% ============================================================



\begin{oefening}[Haakjes maken verschil]{\oefU \ESS}

Vergelijk

\[
A=\neg p\land q
\]

en

\[
B=\neg(p\land q).
\]

Maak voor beide uitdrukkingen een waarheidstabel.

\begin{oplossing}

Voor $A$:

\[
\begin{array}{c|c|c|c}
p&q&\neg p&\neg p\land q\\
\hline
W&W&O&O\\
W&O&O&O\\
O&W&W&W\\
O&O&W&O
\end{array}
\]

Voor $B$:

\[
\begin{array}{c|c|c|c}
p&q&p\land q&\neg(p\land q)\\
\hline
W&W&W&O\\
W&O&O&W\\
O&W&O&W\\
O&O&O&W
\end{array}
\]

De eindkolommen zijn verschillend.

De twee uitdrukkingen hebben dus niet dezelfde logische betekenis.

\end{oplossing}

\end{oefening}





\begin{oefening}[De hoofdoperator herkennen]{\oefU \ESS}

Bepaal voor elke samengestelde uitspraak:

\begin{enumerate}

  \item de elementaire uitspraken;

  \item de gebruikte logische operatoren;

  \item de hoofdoperator.

\end{enumerate}



\begin{question}

\[
\neg p\lor q
\]

\begin{oplossing}

De elementaire uitspraken zijn $p$ en $q$.

De operatoren zijn $\neg$ en $\lor$.

De hoofdoperator is $\lor$, want de volledige uitspraak verbindt

$\neg p$ en $q$ door een disjunctie.

\end{oplossing}

\end{question}



\begin{question}

\[
\neg(p\lor q)
\]

\begin{oplossing}

De elementaire uitspraken zijn $p$ en $q$.

De operatoren zijn $\lor$ en $\neg$.

De hoofdoperator is $\neg$, want de volledige disjunctie wordt ontkend.

\end{oplossing}

\end{question}



\begin{question}

\[
(p\land q)\lor r
\]

\begin{oplossing}

De elementaire uitspraken zijn $p$, $q$ en $r$.

De operatoren zijn $\land$ en $\lor$.

De hoofdoperator is $\lor$.

\end{oplossing}

\end{question}



\begin{question}

\[
p\land(\neg q\lor r)
\]

\begin{oplossing}

De elementaire uitspraken zijn $p$, $q$ en $r$.

De operatoren zijn $\land$, $\neg$ en $\lor$.

De hoofdoperator is $\land$.

\end{oplossing}

\end{question}

\end{oefening}





\begin{oefening}[Vind de fout]{\oefU \ESS}

Een leerling moet de negatie bepalen van

\[
x>3.
\]

Hij schrijft

\[
\neg(x>3)
\quad\Longrightarrow\quad
x<3.
\]

\begin{question}

Is dit correct?

\begin{multipleChoice}

  \choice{Ja.}

  \choice[correct]{Nee.}

\end{multipleChoice}

\end{question}



\begin{question}

Welke waarde toont onmiddellijk dat de negatie onvolledig is?

\begin{oplossing}

\[
x=3.
\]

De oorspronkelijke uitspraak $x>3$ is voor $x=3$ onwaar.

De negatie moet dus voor $x=3$ waar zijn.

Maar $3<3$ is onwaar.

\end{oplossing}

\end{question}



\begin{question}

Geef de correcte negatie.

\begin{oplossing}

\[
x\leq3.
\]

\end{oplossing}

\end{question}



\begin{question}

Welke algemene les volgt hieruit?

\begin{oplossing}

Een uitspraak negeren betekent de volledige verzameling situaties

beschrijven waarin de oorspronkelijke uitspraak niet waar is.

Bij ongelijkheden moet daarom ook goed op de grenswaarde worden gelet.

\end{oplossing}

\end{question}

\end{oefening}





\begin{oefening}[Waar voor welke waarden?]{\oefU \ESS}

Neem $x\in\mathbb{R}$.

Beschouw

\[
p:\quad x>-2
\]

en

\[
q:\quad x\leq4.
\]

\begin{question}

Voor welke waarden van $x$ is

\[
p\land q
\]

waar?

\begin{oplossing}

Beide voorwaarden moeten tegelijk gelden:

\[
x>-2
\]

en

\[
x\leq4.
\]

Dus

\[
-2<x\leq4.
\]

In intervalnotatie:

\[
]-2,4].
\]

\end{oplossing}

\end{question}



\begin{question}

Voor welke waarden van $x$ is

\[
\neg p
\]

waar?

\begin{oplossing}

\[
x\leq-2.
\]

\end{oplossing}

\end{question}



\begin{question}

Voor welke waarden van $x$ is

\[
\neg q
\]

waar?

\begin{oplossing}

\[
x>4.
\]

\end{oplossing}

\end{question}



\begin{question}

Voor welke waarden van $x$ is

\[
\neg p\lor\neg q
\]

waar?

\begin{oplossing}

We hebben

\[
\neg p:\quad x\leq-2
\]

en

\[
\neg q:\quad x>4.
\]

De disjunctie is dus waar voor

\[
x\leq-2
\qquad\text{of}\qquad
x>4.
\]

\end{oplossing}

\end{question}

\end{oefening}





\begin{oefening}[Analyseer de toegangsvoorwaarden]{\oefU \ESS}

Een laboratorium gebruikt een elektronisch toegangscontrolesysteem.

Definieer

\[
p:\quad \text{de badge is geldig}
\]

en

\[
q:\quad \text{de pincode is correct}.
\]

De deur gaat open als

\[
p\land q
\]

waar is.

\begin{question}

Wat gebeurt er wanneer de badge geldig is maar de pincode fout?

\begin{oplossing}

Dan is

\[
p=W,
\qquad
q=O.
\]

Dus

\[
p\land q=O.
\]

De deur blijft gesloten.

\end{oplossing}

\end{question}



\begin{question}

Wat gebeurt er wanneer de badge ongeldig is maar de pincode correct?

\begin{oplossing}

\[
p=O,
\qquad
q=W.
\]

Dus opnieuw

\[
p\land q=O.
\]

De deur blijft gesloten.

\end{oplossing}

\end{question}



\begin{question}

Een programmeur vervangt per ongeluk de voorwaarde door

\[
p\lor q.
\]

Waarom veroorzaakt dit een ernstig beveiligingsprobleem?

\begin{oplossing}

Bij een disjunctie volstaat het dat één voorwaarde waar is.

Daardoor zou iemand toegang kunnen krijgen met alleen een geldige

badge of alleen een correcte pincode.

Het oorspronkelijke systeem eiste bewust dat beide voorwaarden

tegelijk voldaan waren.

\end{oplossing}

\end{question}



\begin{question}

Welke logische operator past dus bij het woord ``beide''?

\begin{oplossing}

De conjunctie:

\[
\land.
\]

\end{oplossing}

\end{question}

\end{oefening}





\begin{oefening}[Een alarmsysteem]{\oefU \ESS}

In een chemische installatie wordt een alarm geactiveerd wanneer

\begin{itemize}

  \item de temperatuur hoger is dan $90\,{}^\circ\mathrm C$;

  \item of de druk hoger is dan $6\,\mathrm{bar}$.

\end{itemize}

Definieer

\[
p:\quad T>90
\]

en

\[
q:\quad P>6.
\]

\begin{question}

Schrijf de alarmvoorwaarde symbolisch.

\begin{oplossing}

\[
p\lor q.
\]

\end{oplossing}

\end{question}



\begin{question}

Gaat het alarm af bij

\[
T=95\,{}^\circ\mathrm C,
\qquad
P=5\,\mathrm{bar}?
\]

\begin{oplossing}

$p$ is waar en $q$ is onwaar.

Dus

\[
p\lor q=W.
\]

Het alarm gaat af.

\end{oplossing}

\end{question}



\begin{question}

Gaat het alarm af bij

\[
T=85\,{}^\circ\mathrm C,
\qquad
P=7\,\mathrm{bar}?
\]

\begin{oplossing}

$p$ is onwaar en $q$ is waar.

Dus

\[
p\lor q=W.
\]

Het alarm gaat af.

\end{oplossing}

\end{question}



\begin{question}

Gaat het alarm af wanneer zowel de temperatuur als de druk te hoog zijn?

\begin{oplossing}

Ja.

Dan zijn $p$ en $q$ beide waar en dus is ook

\[
p\lor q
\]

waar.

Dit illustreert het inclusieve karakter van het wiskundige ``of''.

\end{oplossing}

\end{question}

\end{oefening}





\begin{oefening}[Wat zegt deze voorwaarde precies?]{\oefU \ESS}

Een computerprogramma selecteert meetwaarden waarvoor

\[
x\geq5\land x<12.
\]

\begin{question}

Welke van de volgende waarden worden geselecteerd?

\[
4,\quad5,\quad8,\quad12,\quad15
\]

\begin{oplossing}

De geselecteerde waarden zijn

\[
5\quad\text{en}\quad8.
\]

De waarde $12$ wordt niet geselecteerd omdat de tweede voorwaarde

strikt $x<12$ is.

\end{oplossing}

\end{question}



\begin{question}

Schrijf de selectievoorwaarde in intervalnotatie.

\begin{oplossing}

\[
[5,12[.
\]

\end{oplossing}

\end{question}



\begin{question}

Een programmeur vervangt $\land$ door $\lor$:

\[
x\geq5\lor x<12.
\]

Onderzoek wat er nu gebeurt.

\begin{oplossing}

Voor elk reëel getal is minstens één van beide voorwaarden waar.

Als $x<5$, dan geldt zeker $x<12$.

Als $x\geq5$, dan is de eerste voorwaarde waar.

Dus elk reëel getal voldoet aan

\[
x\geq5\lor x<12.
\]

De nieuwe selectievoorwaarde selecteert dus alle reële getallen.

\end{oplossing}

\end{question}

\end{oefening}





\begin{oefening}[Twee uitdrukkingen vergelijken]{\oefU \ESS}

Onderzoek met een waarheidstabel de uitdrukkingen

\[
A=\neg(p\land q)
\]

en

\[
B=\neg p\lor\neg q.
\]

\begin{question}

Maak de volledige waarheidstabel.

\begin{oplossing}

\[
\begin{array}{c|c|c|c|c|c|c}
p&q&p\land q&\neg(p\land q)&\neg p&\neg q&
\neg p\lor\neg q\\
\hline
W&W&W&O&O&O&O\\
W&O&O&W&O&W&W\\
O&W&O&W&W&O&W\\
O&O&O&W&W&W&W
\end{array}
\]

\end{oplossing}

\end{question}



\begin{question}

Wat valt op wanneer je de laatste relevante kolommen vergelijkt?

\begin{oplossing}

De kolommen van

\[
\neg(p\land q)
\]

en

\[
\neg p\lor\neg q
\]

zijn identiek.

Beide uitdrukkingen hebben dus voor elke mogelijke combinatie van

waarheidswaarden dezelfde waarheidswaarde.

\end{oplossing}

\end{question}



\begin{question}

Waarom is deze conclusie sterker dan wanneer we slechts één keuze

voor $p$ en $q$ zouden controleren?

\begin{oplossing}

De waarheidstabel onderzoekt alle mogelijke combinaties van

waarheidswaarden van $p$ en $q$.

We baseren de conclusie dus niet op één voorbeeld, maar op alle

mogelijke gevallen.

\end{oplossing}

\end{question}

\end{oefening}





% ============================================================

% GEVORDERD

% ============================================================



\begin{oefening}[Een complexe waarheidstabel]{\oefG}

Maak de volledige waarheidstabel van

\[
\neg p\lor(q\land r).
\]

Werk systematisch met tussenkolommen.

\begin{oplossing}

We bepalen eerst $\neg p$ en $q\land r$.

\[
\begin{array}{c|c|c|c|c|c}
p&q&r&\neg p&q\land r&\neg p\lor(q\land r)\\
\hline
W&W&W&O&W&W\\
W&W&O&O&O&O\\
W&O&W&O&O&O\\
W&O&O&O&O&O\\
O&W&W&W&W&W\\
O&W&O&W&O&W\\
O&O&W&W&O&W\\
O&O&O&W&O&W
\end{array}
\]

De eindkolom wordt dus

\[
W,\ O,\ O,\ O,\ W,\ W,\ W,\ W.
\]

\end{oplossing}

\end{oefening}





\begin{oefening}[Ontwerp een logisch regelsysteem]{\oefG}

Een autonome kas beschikt over drie sensoren:

\[
p:\quad \text{de temperatuur is te hoog},
\]

\[
q:\quad \text{de luchtvochtigheid is te laag},
\]

\[
r:\quad \text{het regent buiten}.
\]

Ontwerp een logische voorwaarde voor elk van de volgende acties.

\begin{question}

De ventilatie wordt ingeschakeld wanneer de temperatuur te hoog is

of de luchtvochtigheid te laag is.

\begin{oplossing}

\[
p\lor q.
\]

\end{oplossing}

\end{question}



\begin{question}

Het irrigatiesysteem wordt ingeschakeld wanneer de luchtvochtigheid

te laag is en het buiten niet regent.

\begin{oplossing}

\[
q\land\neg r.
\]

\end{oplossing}

\end{question}



\begin{question}

Het waarschuwingssysteem wordt geactiveerd wanneer de temperatuur

te hoog is en tegelijk minstens één van de twee andere toestanden

optreedt.

\begin{oplossing}

\[
p\land(q\lor r).
\]

\end{oplossing}

\end{question}



\begin{question}

Geef voor de vorige voorwaarde één combinatie van waarheidswaarden

waarbij het waarschuwingssysteem niet geactiveerd wordt, hoewel twee

van de drie elementaire uitspraken waar zijn.

\begin{oplossing}

Neem bijvoorbeeld

\[
p=O,\qquad q=W,\qquad r=W.
\]

Dan is

\[
q\lor r=W,
\]

maar

\[
p\land(q\lor r)
=
O\land W
=
O.
\]

Het systeem wordt dus niet geactiveerd.

Dit toont dat niet alleen het aantal ware uitspraken telt, maar vooral

de logische structuur waarin ze gecombineerd worden.

\end{oplossing}

\end{question}

\end{oefening}





\begin{oefening}[Kan je de fout vinden zonder alles uit te rekenen?]{\oefG}

Een leerling beweert:

\[
\neg(p\lor q)
\]

en

\[
\neg p\lor\neg q
\]

hebben altijd dezelfde waarheidswaarde.

\begin{question}

Zoek één keuze voor $p$ en $q$ waarmee je de bewering kunt weerleggen.

\begin{oplossing}

Neem bijvoorbeeld

\[
p=W,
\qquad
q=O.
\]

Dan is

\[
p\lor q=W,
\]

dus

\[
\neg(p\lor q)=O.
\]

Maar

\[
\neg p=O
\]

en

\[
\neg q=W,
\]

zodat

\[
\neg p\lor\neg q=W.
\]

De twee uitdrukkingen hebben dus verschillende waarheidswaarden.

\end{oplossing}

\end{question}



\begin{question}

Waarom volstaat deze ene combinatie om de bewering ``altijd dezelfde

waarheidswaarde'' te ontkrachten?

\begin{oplossing}

De bewering is algemeen: ze zegt dat de uitdrukkingen voor elke

mogelijke combinatie dezelfde waarheidswaarde hebben.

Eén combinatie waarvoor de waarheidswaarden verschillen is daarom

een tegenvoorbeeld en volstaat om de algemene bewering te weerleggen.

\end{oplossing}

\end{question}



\begin{question}

Wat zou je moeten doen om wél aan te tonen dat twee logische

uitdrukkingen altijd dezelfde waarheidswaarde hebben?

\begin{oplossing}

We moeten alle mogelijke combinaties van waarheidswaarden onderzoeken.

Bij een eindig aantal elementaire uitspraken kan dat systematisch met

een waarheidstabel.

\end{oplossing}

\end{question}

\end{oefening}





\begin{oefening}[Analyse van een medische beslisregel]{\oefG}

Een automatisch bewakingssysteem op een intensieve-zorgafdeling

registreert twee vereenvoudigde voorwaarden:

\[
p:\quad \text{de zuurstofsaturatie ligt onder de ingestelde grens},
\]

\[
q:\quad \text{de hartfrequentie ligt buiten het ingestelde bereik}.
\]

Dit is uitsluitend een wiskundig model van de logische structuur van

een bewakingssysteem.

Een eerste ontwerp activeert een waarschuwing wanneer

\[
p\lor q.
\]

\begin{question}

Beschrijf in woorden wanneer de waarschuwing wordt geactiveerd.

\begin{oplossing}

De waarschuwing wordt geactiveerd wanneer minstens één van beide

voorwaarden optreedt:

\begin{itemize}

  \item de zuurstofsaturatie ligt onder de ingestelde grens;

  \item de hartfrequentie ligt buiten het ingestelde bereik;

  \item of beide situaties treden tegelijk op.

\end{itemize}

\end{oplossing}

\end{question}



\begin{question}

Een tweede ontwerp gebruikt

\[
p\land q.
\]

Wat is het fundamentele verschil?

\begin{oplossing}

Bij

\[
p\land q
\]

wordt alleen gewaarschuwd als beide voorwaarden tegelijk waar zijn.

Dit is dus een veel strengere voorwaarde dan

\[
p\lor q,
\]

waar één afwijkende toestand al voldoende is.

\end{oplossing}

\end{question}



\begin{question}

Waarom kan een kleine verandering van een logische operator in een

technologisch systeem grote gevolgen hebben?

\begin{oplossing}

De operator bepaalt welke combinaties van toestanden tot een bepaalde

actie leiden.

Het vervangen van $\lor$ door $\land$, of omgekeerd, verandert dus de

volledige beslisregel en kan ertoe leiden dat een systeem te vaak of

juist te weinig reageert.

\end{oplossing}

\end{question}

\end{oefening}





\begin{oefening}[Construeer zelf]{\oefG}

Maak zelf drie elementaire wiskundige uitspraken $p$, $q$ en $r$.

Construeer vervolgens een samengestelde uitspraak waarin minstens

\[
\neg,\qquad\land,\qquad\lor
\]

elk minstens één keer voorkomen.

\begin{question}

Noteer de drie elementaire uitspraken en geef hun waarheidswaarden.

\end{question}



\begin{question}

Schrijf je samengestelde uitspraak symbolisch.

\end{question}



\begin{question}

Vertaal de samengestelde uitspraak volledig naar woorden.

\end{question}



\begin{question}

Bepaal de waarheidswaarde van de samengestelde uitspraak en motiveer

stap voor stap.

\end{question}



\begin{oplossing}

Er zijn veel correcte antwoorden mogelijk.

Een voorbeeld:

\[
p:\quad 8\text{ is even}\qquad(W),
\]

\[
q:\quad 8\text{ is een priemgetal}\qquad(O),
\]

\[
r:\quad 8>5\qquad(W).
\]

Neem

\[
\neg q\land(p\lor r).
\]

We hebben

\[
\neg q=W
\]

en

\[
p\lor r=W.
\]

Daarom:

\[
\neg q\land(p\lor r)
=
W\land W
=
W.
\]

Een correct eigen antwoord moet dezelfde logische stappen duidelijk

verantwoorden.

\end{oplossing}

\end{oefening}





% ============================================================

% SYNTHESE

% ============================================================



\begin{oefening}[Synthese: van taal tot logische analyse]{\oefU \ESS}

Beschouw voor $x\in\mathbb{R}$ de voorwaarden

\[
p:\quad x\geq-1
\]

en

\[
q:\quad x<4.
\]

\begin{question}

Vertaal

\[
p\land q
\]

naar een voorwaarde voor $x$.

\begin{oplossing}

\[
x\geq-1\land x<4.
\]

Dus

\[
-1\leq x<4.
\]

\end{oplossing}

\end{question}



\begin{question}

Welke van de waarden

\[
-2,\quad-1,\quad0,\quad4,\quad5
\]

maken $p\land q$ waar?

\begin{oplossing}

Alleen

\[
-1\quad\text{en}\quad0.
\]

\end{oplossing}

\end{question}



\begin{question}

Geef de negatie van $p$.

\begin{oplossing}

\[
\neg p:\quad x<-1.
\]

\end{oplossing}

\end{question}



\begin{question}

Geef de negatie van $q$.

\begin{oplossing}

\[
\neg q:\quad x\geq4.
\]

\end{oplossing}

\end{question}



\begin{question}

Beschrijf wanneer

\[
\neg p\lor\neg q
\]

waar is.

\begin{oplossing}

De uitspraak is waar wanneer

\[
x<-1
\]

of

\[
x\geq4.
\]

Dus voor

\[
x\in]-\infty,-1[\ \cup\ [4,+\infty[.
\]

\end{oplossing}

\end{question}



\begin{question}

Vergelijk dit antwoord met het gebied waarin $p\land q$ waar is.

Wat merk je op?

\begin{oplossing}

$p\land q$ is waar voor

\[
-1\leq x<4.
\]

$\neg p\lor\neg q$ is precies waar buiten dat interval.

We zien hier dus opnieuw het patroon dat later formeel wordt beschreven

door een wet van De Morgan:

\[
\neg(p\land q)
\]

heeft dezelfde waarheidswaarden als

\[
\neg p\lor\neg q.
\]

\end{oplossing}

\end{question}

\end{oefening}





\end{document}

